% math0 manuscript. Build with: latexmk -pdf paper.tex
\documentclass[11pt]{amsart}
\usepackage[margin=1.1in]{geometry}
\usepackage{amsmath,amssymb,amsthm,mathtools}
\usepackage{url}
\usepackage{hyperref}
\usepackage{xcolor}
\usepackage{listings}
\usepackage{newunicodechar}

\lstdefinelanguage{lean}{
  keywords={theorem,lemma,def,abbrev,example,fun,by,have,show,from,let,in,if,then,else,match,with,
    namespace,end,open,import,section,variable,instance,structure,where,calc,obtain,intro,exact,
    refine,rcases,rintro,simp,omega,decide,norm_num,linarith,nlinarith,positivity,induction,cases,
    private},
  sensitive=true, comment=[l]{--}, morecomment=[s]{/-}{-/}, string=[b]"}
\lstset{language=lean, basicstyle=\ttfamily\scriptsize, columns=fullflexible, keepspaces=true, breaklines=true,
  keywordstyle=\color{blue!60!black}\bfseries, commentstyle=\color{gray}, frame=single,
  extendedchars=true, inputencoding=utf8,
  literate={ℕ}{{$\mathbb{N}$}}1 {ℤ}{{$\mathbb{Z}$}}1 {∀}{{$\forall$}}1 {∃}{{$\exists$}}1
    {→}{{$\to$}}1 {∧}{{$\wedge$}}1 {≤}{{$\le$}}1 {∣}{{$\mid$}}1 {∑}{{$\sum$}}1
    {∈}{{$\in$}}1 {×}{{$\times$}}1}
% Unicode commonly found in Lean sources.
\newunicodechar{ℕ}{\ensuremath{\mathbb{N}}}
\newunicodechar{ℤ}{\ensuremath{\mathbb{Z}}}
\newunicodechar{ℚ}{\ensuremath{\mathbb{Q}}}
\newunicodechar{ℝ}{\ensuremath{\mathbb{R}}}
\newunicodechar{ℂ}{\ensuremath{\mathbb{C}}}
\newunicodechar{∀}{\ensuremath{\forall}}
\newunicodechar{∃}{\ensuremath{\exists}}
\newunicodechar{→}{\ensuremath{\to}}
\newunicodechar{←}{\ensuremath{\leftarrow}}
\newunicodechar{↔}{\ensuremath{\leftrightarrow}}
\newunicodechar{↦}{\ensuremath{\mapsto}}
\newunicodechar{¬}{\ensuremath{\neg}}
\newunicodechar{∧}{\ensuremath{\wedge}}
\newunicodechar{∨}{\ensuremath{\vee}}
\newunicodechar{≤}{\ensuremath{\le}}
\newunicodechar{≥}{\ensuremath{\ge}}
\newunicodechar{≠}{\ensuremath{\ne}}
\newunicodechar{∈}{\ensuremath{\in}}
\newunicodechar{∉}{\ensuremath{\notin}}
\newunicodechar{⊆}{\ensuremath{\subseteq}}
\newunicodechar{∩}{\ensuremath{\cap}}
\newunicodechar{∪}{\ensuremath{\cup}}
\newunicodechar{∑}{\ensuremath{\sum}}
\newunicodechar{∏}{\ensuremath{\prod}}
\newunicodechar{⟨}{\ensuremath{\langle}}
\newunicodechar{⟩}{\ensuremath{\rangle}}
\newunicodechar{·}{\ensuremath{\cdot}}
\newunicodechar{∘}{\ensuremath{\circ}}
\newunicodechar{⁻}{\ensuremath{^{-}}}
\newunicodechar{¹}{\ensuremath{^{1}}}
\newunicodechar{²}{\ensuremath{^{2}}}
\newunicodechar{₀}{\ensuremath{_{0}}}
\newunicodechar{₁}{\ensuremath{_{1}}}
\newunicodechar{₂}{\ensuremath{_{2}}}
\newunicodechar{α}{\ensuremath{\alpha}}
\newunicodechar{β}{\ensuremath{\beta}}
\newunicodechar{ε}{\ensuremath{\varepsilon}}
\newunicodechar{δ}{\ensuremath{\delta}}
\newunicodechar{π}{\ensuremath{\pi}}
\newunicodechar{σ}{\ensuremath{\sigma}}
\newunicodechar{φ}{\ensuremath{\varphi}}
\newunicodechar{∞}{\ensuremath{\infty}}
\newunicodechar{⌊}{\ensuremath{\lfloor}}
\newunicodechar{⌋}{\ensuremath{\rfloor}}
\newunicodechar{⌈}{\ensuremath{\lceil}}
\newunicodechar{⌉}{\ensuremath{\rceil}}
\newunicodechar{∣}{\ensuremath{\mid}}
\newunicodechar{≡}{\ensuremath{\equiv}}
\newunicodechar{×}{\ensuremath{\times}}
\newunicodechar{⧸}{\ensuremath{/}}
\newunicodechar{‖}{\ensuremath{\|}}
\newunicodechar{√}{\ensuremath{\surd}}
\newunicodechar{∅}{\ensuremath{\emptyset}}
\newunicodechar{⦃}{\ensuremath{\{\!|}}
\newunicodechar{⦄}{\ensuremath{|\!\}}}
\newunicodechar{𝓝}{\ensuremath{\mathcal{N}}}
\newunicodechar{≫}{\ensuremath{\gg}}
\newunicodechar{≪}{\ensuremath{\ll}}

\theoremstyle{plain}
\newtheorem{theorem}{Theorem}
\newtheorem{lemma}[theorem]{Lemma}
\newtheorem{proposition}[theorem]{Proposition}
\newtheorem{corollary}[theorem]{Corollary}
\theoremstyle{definition}
\newtheorem{definition}[theorem]{Definition}
\newtheorem{remark}[theorem]{Remark}

\emergencystretch=3em
% Lean identifiers in running text (breakable at dots and underscores).
\expandafter\def\expandafter\UrlBreaks\expandafter{\UrlBreaks\do\_\do\.}
\DeclareUrlCommand\decl{\urlstyle{tt}}
\newcommand{\N}{\mathbb{N}}
\newcommand{\Z}{\mathbb{Z}}
\newcommand{\cp}{\operatorname{cp}}

\title[The affine 42-block structure of OEIS A351743]{From $n=139$ on, OEIS A351743 is affine in blocks of length $42$: proof of the conjectures of Sycamore and Sigrist}
\author{math0}
\address{Produced by Claude (an AI model developed by Anthropic) in the math0 project}
\date{October 20, 2026}

\begin{document}
\begin{abstract}
Let $a(1)=1$ and let $a(n+1)$ be the greatest divisor of $s(n)=a(1)+\dots+a(n)$ that is prime to $a(n)$; this is OEIS A351743, and $s$ is A355149.
We prove that there are explicit natural numbers $P_i,Q_i,R_i,U_i$ ($0\le i<42$) with $a(139+42m+i)=P_iY_m+Q_i$ and $s(139+42m+i)=R_iY_m+U_i$ for all $m\ge0$, where $Y_m=(3645^m-1)/3644$.
This proves the four conjectures recorded in the two entries in 2022: $1$ and $2$ occur infinitely often (D.~J. Sycamore); $a(2(n+21)-1)=a(2n-1)$ and $a(2(n+21))=3645\,a(2n)+b(n)$ with $b$ $21$-periodic, for $n\ge70$ (R.~Sigrist); and $s(n+42)=3645\,s(n)+c(n)$ with $c$ $42$-periodic, for $n\ge138$ (A355149).
To our knowledge the result is new; both entries still list these as conjectures.
The proof is verified in Lean~4 against a frozen, fidelity-reviewed statement that math0 formalized from the OEIS entries (not part of the formal-conjectures repository).
\end{abstract}
\maketitle

\section{Introduction}

The sequence A351743 of the OEIS \cite{A351743}, contributed by David James Sycamore on June~7, 2022, is defined by
\[
  a(1)=1,\qquad a(n+1)=\text{the greatest divisor of } s(n) \text{ that is prime to } a(n)\quad(n\ge1),
\]
where $s(n)=a(1)+\dots+a(n)$. It begins $1,1,2,1,5,2,3,5,4,3,1,28,1,57,2,\dots$, and the partial sums $s(n)$ form A355149 \cite{A355149} (M.~De Vlieger, R.~Sigrist and D.~J. Sycamore, June~21, 2022).
Equivalently, $a(n+1)$ is $s(n)$ with every prime power $p^k\,\|\,s(n)$, $p\mid a(n)$, removed; as $a(n)\mid s(n)$, $a(n+1)=1$ exactly when $s(n)$ and $a(n)$ have the same prime divisors.
Sycamore conjectured that $1$, and hence $2$, occur infinitely often in $a$; R.~Sigrist conjectured on June~20, 2022 that $a(2(n+21)-1)=a(2n-1)$ and $a(2(n+21))=3645\,a(2n)+b(n)$ with $b$ $21$-periodic, for $n\ge70$; and A355149 records the conjecture $s(n+42)=3645\,s(n)+b(n)$ with $b$ $42$-periodic, for $n\ge138$ (the exact wording is quoted in Section~\ref{sec:statement}).
We prove all four (Theorem~\ref{thm:main}). They follow from an exact description of the sequence from $n=139$ on (Theorem~\ref{thm:block}): on the $m$-th block of $42$ consecutive indices, $a$ and $s$ are affine functions of $Y_m=(3645^m-1)/3644$.
In particular, for odd $n\ge139$ the terms $a(n)$ are $42$-periodic with values in $\{1,\dots,5\}$ (eight $1$'s and five $2$'s per period), while for even $n$ they grow by a factor of about $3645=3^6\cdot5$ per block.

\subsection*{Method}
We exhibit $42$ quadruples $(P_i,Q_i,R_i,U_i)$ (Table~\ref{tab:rows}) such that the recursion $(a(n),s(n))\mapsto(a(n+1),s(n+1))$ maps $(P_iW+Q_i,\,R_iW+U_i)$ to the pair of the next quadruple for \emph{every} $W\in\N$, and the last one to the first one at $3645W+1$. Each step follows from Lemma~\ref{lem:cp} and finitely many facts about explicit integers, checked by the Lean kernel; the first $139$ terms are computed directly. As the steps hold for all $W$, no $p$-adic information on $3645^m-1$ is needed.

\subsection*{Status and prior work}
As of October~2026, A351743 (revision \#43, June~20, 2022) and A355149 (revision \#11, November~17, 2022) still list all four statements as conjectures. Apart from A355149, the only OEIS entry citing A351743 is A357735 \cite{A357735}, which says nothing about them.
The math0 referees found no mention of A351743, A355149, $3645$ or Sycamore in the full text of S.~Fried's papers proving OEIS conjectures \cite{Fr24,Fr25,Fr26,Fr26b}, and no proof or counterexample on the web, arXiv, MathOverflow, Mathematics Stack Exchange or the SeqFan list; the statements are not in the formal-conjectures repository \cite{FC}.
We therefore classify the result as \emph{new}; the structure we prove is the one Sigrist observed numerically, and our contribution is a (computer-assisted) proof that it persists.

\section{Statement}\label{sec:statement}

\subsection*{The conjectures}
The A351743 entry \cite{A351743} contains the following lines, verbatim (we omit the A-number after each line tag):
\begin{quote}\raggedright\frenchspacing
{\footnotesize\texttt{\%N a(1)=1. Thereafter, a(n+1) is the greatest divisor of s(n) which is prime to a(n), where s(n) is the n-th partial sum.}\\
\texttt{\%C Conjecture: a(n)=1 occurs infinitely many times, and therefore so does a(n)=2.}\\
\texttt{\%F Conjectures from \_R\'emy Sigrist\_, Jun 20 2022: (Start)}\\
\texttt{\%F - a(2*(n+21)-1) = a(2*n-1) for n >= 70,}\\
\texttt{\%F - a(2*(n+21)) = 3645*a(2*n) + b(n) for n >= 70 (where b is 21-periodic).}\\
\texttt{\%F (End)}}
\end{quote}
The entry is by \texttt{\_David James Sycamore\_, Jun 07 2022}; the COMMENTS conjecture is unsigned, so we attribute it to him. The A355149 entry \cite{A355149}, ``Partial sums of A351743'', with offset $1$, contains:
\begin{quote}\raggedright\frenchspacing
{\footnotesize\texttt{\%F Conjecture: a(n + 42) = 3645*a(n) + b(n) for n >= 138 (where b is 42-periodic).}}
\end{quote}
Here \texttt{a} is our $s$; we rename this $b$ to $c$.

\begin{theorem}\label{thm:main}
Let $a$ and $s$ be as in the introduction. Then:
\begin{enumerate}
\item[(i)] $a(n)=1$ for infinitely many $n$, and $a(n)=2$ for infinitely many $n$;
\item[(ii)] $a(2(n+21)-1)=a(2n-1)$ for every $n\ge70$;
\item[(iii)] $a(2(n+21))=3645\,a(2n)+b(n)$ for all $n\ge70$, for some $b\colon\N\to\Z$ with $b(n+21)=b(n)$;
\item[(iv)] $s(n+42)=3645\,s(n)+c(n)$ for all $n\ge138$, for some $c\colon\N\to\Z$ with $c(n+42)=c(n)$.
\end{enumerate}
\end{theorem}

\subsection*{Lean formalization}
The statement was formalized by math0 in \path{lean/Math0/Statements/OEIS/A351743.lean}, reviewed for fidelity and frozen (SHA-256 \texttt{41c82d15\ldots 47d3}) before any proof attempt. Its definitions, verbatim (import, docstrings, comments, blank lines and theorems omitted):
\begin{lstlisting}
namespace Math0.Statements.A351743
def strip : ℕ → ℕ → ℕ → ℕ
  | 0, w, _ => w
  | fuel + 1, w, v => if Nat.gcd w v = 1 then w else strip fuel (w / Nat.gcd w v) v
def coprimePart (w v : ℕ) : ℕ := strip w w v
def step (p : ℕ × ℕ) : ℕ × ℕ :=
  match p with
  | (x, y) =>
    match coprimePart y x with
    | z => (z, y + z)
def state (n : ℕ) : ℕ × ℕ := step^[n] (1, 1)
def a : ℕ → ℕ
  | 0 => 0
  | n + 1 => (state n).1
def s (n : ℕ) : ℕ := ∑ k ∈ Finset.Icc 1 n, a k
def sycamore : Prop :=
  {n : ℕ | a n = 1}.Infinite ∧ {n : ℕ | a n = 2}.Infinite
def sigristOdd : Prop :=
  ∀ n : ℕ, 70 ≤ n → a (2 * (n + 21) - 1) = a (2 * n - 1)
def sigristEven : Prop :=
  ∃ b : ℕ → ℤ, (∀ n, b (n + 21) = b n) ∧
    ∀ n : ℕ, 70 ≤ n → (a (2 * (n + 21)) : ℤ) = 3645 * (a (2 * n) : ℤ) + b n
def partialSums : Prop :=
  ∃ c : ℕ → ℤ, (∀ n, c (n + 42) = c n) ∧
    ∀ n : ℕ, 138 ≤ n → (s (n + 42) : ℤ) = 3645 * (s n : ℤ) + c n
def statement : Prop :=
  sycamore ∧ sigristOdd ∧ sigristEven ∧ partialSums
end Math0.Statements.A351743
\end{lstlisting}
Here \decl{a} has the OEIS offset $1$; \texttt{a 0 = 0} is a junk value that no conjecture reaches (and it is neither $1$ nor $2$), and \texttt{s n} is literally $\sum_{k=1}^n a(k)$. As in the entry's PARI program, \texttt{coprimePart w v} divides $w$ by $\gcd(w,v)$ until this gcd is $1$ (the fuel $w$ suffices, as each division at least halves $w$), and \texttt{state n} is $(a(n+1),s(n+1))$. The file proves, among others (proofs omitted):
\begin{lstlisting}
theorem coprimePart_isGreatest (w v : ℕ) (hw : 0 < w) :
    IsGreatest {d | d ∣ w ∧ Nat.Coprime d v} (coprimePart w v)
theorem a_one : a 1 = 1
theorem a_succ_isGreatest (n : ℕ) (hn : 1 ≤ n) :
    IsGreatest {d | d ∣ s n ∧ Nat.Coprime d (a n)} (a (n + 1))
\end{lstlisting}
The theorems \decl{a_one} and \decl{a_succ_isGreatest} are the OEIS definition word for word and determine $a(n)$ for all $n\ge1$. Three theorems proved by \texttt{decide +kernel} compare with the OEIS data: \decl{a_data} checks all $67$ listed terms of A351743, \decl{s_data} all $50$ listed terms of A355149, and \decl{a_example} the values $a(150),\dots,a(153)=473394,1,2840365,2$ of the entry's last example.
The four conjectures are encoded literally. In \decl{sigristOdd} the truncated subtraction is harmless, since $2n-1\ge139$. The periodic sequences are integer-valued, since their signs are not specified, and the relations are stated in $\Z$. ``$b$ is $21$-periodic'' is read as $b(n+21)=b(n)$ for all $n\in\N$; this is equivalent to periodicity for $n\ge70$, since only the values at $n\ge70$ enter and they extend backwards periodically. Thus \decl{sigristEven} says exactly that $a(2(n+21))-3645\,a(2n)$ depends only on $n\bmod21$ for $n\ge70$; likewise for $c$ in \decl{partialSums}.

\section{Proof}\label{sec:proof}

The proof follows the Lean file \path{lean/Math0/Frontier/OEIS/A351743.lean}; we give the names of the corresponding Lean declarations.
For $w\ge1$ and $v\ge0$ let $\cp(w,v)$ be the greatest divisor of $w$ prime to $v$ (\decl{coprimePart}), and for $(x,y)\in\N^2$, $y\ge1$, put $T(x,y)=(\cp(y,x),\ y+\cp(y,x))$ (\decl{step}). By definition, $(a(n+1),s(n+1))=T^n(1,1)$ for all $n\ge0$ (\decl{state}, \decl{state_snd}).

\subsection{Two arithmetic lemmas}

\begin{lemma}[\decl{cp_eq}]\label{lem:cp}
Let $d,e\ge1$ and $v,k\ge0$. If $\gcd(d,v)=1$ and $e\mid v^k$, then $\cp(de,v)=d$.
\end{lemma}

\begin{proof}
The number $d$ divides $de$ and is prime to $v$. Conversely, let $x\mid de$ with $\gcd(x,v)=1$. Then $\gcd(x,v^k)=1$, hence $\gcd(x,e)=1$ because $e\mid v^k$, and so $x\mid d$ and $x\le d$. Thus $d$ is the greatest element of $\{x: x\mid de,\ \gcd(x,v)=1\}$, which is $\cp(de,v)$ by \decl{coprimePart_isGreatest}.
\end{proof}

\begin{lemma}[\decl{cop_lin}]\label{lem:lin}
Let $L,M,g,K,k,W\ge0$. If $g\mid L$, $\gcd(M,g)=1$ and $K\mid g^k$, then $\gcd(LW+M,K)=1$.
\end{lemma}

\begin{proof}
Since $LW+M\equiv M\pmod g$, we have $\gcd(LW+M,g)=\gcd(M,g)=1$, hence $\gcd(LW+M,g^k)=1$, and $K\mid g^k$ gives the claim.
\end{proof}

\subsection{Rows and certified steps}
A \emph{row} is a quadruple $r=(P,Q,R,U)\in\N^4$. For $W\in\N$ put $F_r(W)=(PW+Q,\ RW+U)$ (\decl{F}).
For rows $r=(P,Q,R,U)$ and $r'=(P',Q',R',U')$ we consider two kinds of certificates.

\begin{definition}[\decl{GoodSmall}, \decl{GoodBig}]\label{def:good}
\emph{(S)} A \emph{small step} from $r$ to $r'$ with parameters $(e,g,k)$ means:
$P=0$, $R=eP'$, $U=eQ'$, $e\mid Q^k$, $g\mid P'$, $\gcd(Q',g)=1$, $Q\mid g^k$, $e\ge1$, $Q'\ge1$, $R'=R+P'$ and $U'=U+Q'$.

\emph{(B)} A \emph{big step} from $r$ to $r'$ with parameters $(\lambda,f,g,h,k)$ means:
$P'=0$, $R=\lambda P$, $U=\lambda Q$, $\lambda=Q'f$, $h\mid P$, $h\mid Q$, $f\mid h^k$, $g\mid P$, $\gcd(Q,g)=1$, $Q'\mid g^k$, $Q'\ge1$, $Q\ge1$, $f\ge1$, $R'=R$ and $U'=U+Q'$.
\end{definition}

In a small step the current term $a(n)=Q$ is a constant, and $s(n)=e\cdot a(n+1)$ where $e$ is built from the primes of $Q$. In a big step $s(n)=\lambda\,a(n)$ for a constant $\lambda$, and the next term is the constant $Q'$.

\begin{lemma}[\decl{step_small}, \decl{step_big}, \decl{step_good}]\label{lem:step}
If there is a small or a big step from $r$ to $r'$, then $T(F_r(W))=F_{r'}(W)$ for every $W\in\N$.
\end{lemma}

\begin{proof}
\emph{Small step.} Since $P=0$, $F_r(W)=(Q,\ RW+U)$ and $RW+U=e\,(P'W+Q')$. Here $P'W+Q'\ge Q'\ge1$, $e\ge1$ and $e\mid Q^k$, and $\gcd(P'W+Q',Q)=1$ by Lemma~\ref{lem:lin} (with $L=P'$, $M=Q'$, $K=Q$). Lemma~\ref{lem:cp} gives $\cp(RW+U,Q)=P'W+Q'$. The new partial sum is $RW+U+P'W+Q'=(R+P')W+(U+Q')=R'W+U'$.

\emph{Big step.} Put $v=PW+Q\ge Q\ge1$. Then $RW+U=\lambda v=Q'\cdot(fv)$. Since $h\mid P$ and $h\mid Q$ we have $h\mid v$, so $f\mid h^k\mid v^k$ and $fv\mid v^{k+1}$. By Lemma~\ref{lem:lin} (with $L=P$, $M=Q$, $K=Q'$), $\gcd(v,Q')=1$. Lemma~\ref{lem:cp} with $d=Q'$ and $e=fv\ge1$ gives $\cp(RW+U,v)=Q'=P'W+Q'$, since $P'=0$. The new partial sum is $RW+U+Q'=R'W+U'$.
\end{proof}

\subsection{The table}
Let $r_i=(P_i,Q_i,R_i,U_i)$, $0\le i\le41$, be the rows of Table~\ref{tab:rows} (\decl{rows}). For the shift $\sigma(P,Q,R,U)=(3645P,\ P+Q,\ 3645R,\ R+U)$ (\decl{shift}) we have $F_{\sigma(r)}(W)=F_r(3645W+1)$ for all $W$ (\decl{F_shift}). Let $r'_i=r_{i+1}$ for $i<41$ and $r'_{41}=\sigma(r_0)=(0,\,3,\,4191876044100,\,1150350171)$ (\decl{nextRow}).

\begin{lemma}[\decl{good_all}]\label{lem:table}
For every $i<42$ there is a step from $r_i$ to $r'_i$ with the parameters of Table~\ref{tab:rows}: a small step (S) for even $i$ and a big step (B) for odd $i$, in both cases with $k=4$.
\end{lemma}

\begin{proof}
Each condition is an equality, divisibility, gcd value or inequality between explicit natural numbers below $5\cdot10^{12}$ (the powers $Q^4,g^4,h^4$ are at most $15^4$). In Lean, $\forall i<42,\ \mathrm{Good}(r_i,r'_i,\mathrm{cert}_i)$ is proved by \texttt{decide +kernel}: the kernel evaluates the decision procedure, using its built-in arithmetic on numerals.
\end{proof}

\begin{table}[t]
\centering
\scriptsize
\setlength{\tabcolsep}{3.2pt}
\begin{tabular}{rrrrrrcrrrrrrrrr}
\multicolumn{6}{c}{even $i$: $P_i=0$, small step (S)} && \multicolumn{9}{c}{odd $i$: big step (B)}\\
$i$ & $Q_i$ & $R_i$ & $U_i$ & $e$ & $g$ && $i$ & $P_i$ & $Q_i$ & $R_i$ & $U_i$ & $\lambda$ & $f$ & $g$ & $h$\\\hline
0 & 3 & 1150034580 & 315591 & 3 & 3 && 1 & 383344860 & 105197 & 1533379440 & 420788 & 4 & 1 & 2 & 1 \\
2 & 4 & 1533379440 & 420792 & 8 & 2 && 3 & 191672430 & 52599 & 1725051870 & 473391 & 9 & 9 & 1 & 3 \\
4 & 1 & 1725051870 & 473392 & 1 & 1 && 5 & 1725051870 & 473392 & 3450103740 & 946784 & 2 & 2 & 1 & 2 \\
6 & 1 & 3450103740 & 946785 & 1 & 1 && 7 & 3450103740 & 946785 & 6900207480 & 1893570 & 2 & 1 & 2 & 15 \\
8 & 2 & 6900207480 & 1893572 & 4 & 2 && 9 & 1725051870 & 473393 & 8625259350 & 2366965 & 5 & 1 & 5 & 1 \\
10 & 5 & 8625259350 & 2366970 & 5 & 5 && 11 & 1725051870 & 473394 & 10350311220 & 2840364 & 6 & 6 & 1 & 6 \\
12 & 1 & 10350311220 & 2840365 & 1 & 1 && 13 & 10350311220 & 2840365 & 20700622440 & 5680730 & 2 & 1 & 2 & 5 \\
14 & 2 & 20700622440 & 5680732 & 4 & 2 && 15 & 5175155610 & 1420183 & 25875778050 & 7100915 & 5 & 1 & 5 & 1 \\
16 & 5 & 25875778050 & 7100920 & 5 & 5 && 17 & 5175155610 & 1420184 & 31050933660 & 8521104 & 6 & 2 & 3 & 2 \\
18 & 3 & 31050933660 & 8521107 & 3 & 3 && 19 & 10350311220 & 2840369 & 41401244880 & 11361476 & 4 & 1 & 2 & 1 \\
20 & 4 & 41401244880 & 11361480 & 8 & 2 && 21 & 5175155610 & 1420185 & 46576400490 & 12781665 & 9 & 9 & 1 & 15 \\
22 & 1 & 46576400490 & 12781666 & 1 & 1 && 23 & 46576400490 & 12781666 & 93152800980 & 25563332 & 2 & 2 & 1 & 2 \\
24 & 1 & 93152800980 & 25563333 & 1 & 1 && 25 & 93152800980 & 25563333 & 186305601960 & 51126666 & 2 & 1 & 2 & 3 \\
26 & 2 & 186305601960 & 51126668 & 4 & 2 && 27 & 46576400490 & 12781667 & 232882002450 & 63908335 & 5 & 1 & 5 & 1 \\
28 & 5 & 232882002450 & 63908340 & 5 & 5 && 29 & 46576400490 & 12781668 & 279458402940 & 76690008 & 6 & 6 & 1 & 6 \\
30 & 1 & 279458402940 & 76690009 & 1 & 1 && 31 & 279458402940 & 76690009 & 558916805880 & 153380018 & 2 & 1 & 2 & 1 \\
32 & 2 & 558916805880 & 153380020 & 4 & 2 && 33 & 139729201470 & 38345005 & 698646007350 & 191725025 & 5 & 5 & 1 & 5 \\
34 & 1 & 698646007350 & 191725026 & 1 & 1 && 35 & 698646007350 & 191725026 & 1397292014700 & 383450052 & 2 & 2 & 1 & 6 \\
36 & 1 & 1397292014700 & 383450053 & 1 & 1 && 37 & 1397292014700 & 383450053 & 2794584029400 & 766900106 & 2 & 1 & 2 & 1 \\
38 & 2 & 2794584029400 & 766900108 & 4 & 2 && 39 & 698646007350 & 191725027 & 3493230036750 & 958625135 & 5 & 1 & 5 & 1 \\
40 & 5 & 3493230036750 & 958625140 & 5 & 5 && 41 & 698646007350 & 191725028 & 4191876044100 & 1150350168 & 6 & 2 & 3 & 2 \\
\end{tabular}
\caption{The rows $r_i=(P_i,Q_i,R_i,U_i)$, $0\le i\le 41$, with $a(139+42m+i)=P_iY_m+Q_i$ and $s(139+42m+i)=R_iY_m+U_i$ (left: even $i$, where $P_i=0$; right: odd $i$), and the parameters of the step $r_i\to r'_i$ (Definition~\ref{def:good}; $k=4$ throughout). These are the lists \texttt{rows} and \texttt{certs} of the Lean file.}\label{tab:rows}
\end{table}

\subsection{The block theorem}

\begin{theorem}[\decl{key}, \decl{a_form}, \decl{s_form}]\label{thm:block}
Let $Y_0=0$ and $Y_{m+1}=3645\,Y_m+1$, so that $Y_m=(3645^m-1)/3644$. For all $m\ge0$ and $0\le i\le41$,
\[
  a(139+42m+i)=P_iY_m+Q_i,\qquad s(139+42m+i)=R_iY_m+U_i .
\]
\end{theorem}

\begin{proof}
Equivalently, $T^{138+42m+i}(1,1)=F_{r_i}(Y_m)$. We prove this by induction on $m$, and for each $m$ by induction on $i$.

\emph{Base.} $T^{138}(1,1)=(a(139),s(139))=(3,315591)=F_{r_0}(0)$ (\decl{state_base}: \texttt{decide +kernel} evaluates the $138$ steps in the kernel; the partial sums involved are at most $s(139)=315591$).

\emph{From $i$ to $i+1$, for $i<41$.} By the induction hypothesis and Lemmas~\ref{lem:step} and~\ref{lem:table} with $W=Y_m$,
$T^{138+42m+i+1}(1,1)=T(F_{r_i}(Y_m))=F_{r_{i+1}}(Y_m)$.

\emph{From block $m$ to block $m+1$.} Since $138+42(m+1)=(138+42m+41)+1$, Lemmas~\ref{lem:step} and~\ref{lem:table} and \decl{F_shift} give $T^{138+42(m+1)}(1,1)=T(F_{r_{41}}(Y_m))=F_{\sigma(r_0)}(Y_m)=F_{r_0}(3645Y_m+1)=F_{r_0}(Y_{m+1})$.
\end{proof}

By Table~\ref{tab:rows}, $P_i=0$ for every even $i$ (\decl{row_even_const}, checked by \texttt{decide}), with
\[
  (Q_0,Q_2,\dots,Q_{40})=(3,4,1,1,2,5,1,2,5,3,4,1,1,2,5,1,2,1,1,2,5),
\]
whereas $P_i>0$ and $s(n)=\lambda_i\,a(n)$ with $\lambda_i\in\{2,4,5,6,9\}$ for odd $i$.
For instance $a(150)=Q_{11}=473394$, $a(151)=Q_{12}=1$, $a(152)=Q_{13}=2840365$ and $a(153)=Q_{14}=2$ are the values in the last example of the OEIS entry.

\subsection{Proof of Theorem~\ref{thm:main}}

\emph{(i)} (\decl{one_inf}, \decl{two_inf}.) As $r_4=(0,1,\dots)$ and $r_8=(0,2,\dots)$, Theorem~\ref{thm:block} gives $a(143+42m)=1$ and $a(147+42m)=2$ for all $m\ge0$.

\emph{(ii)} (\decl{odd_part}.) Let $n\ge70$. Then $2n-1\ge139$, so $2n-1=139+42m+i$ with $m\ge0$ and $0\le i<42$, and $i$ is even because $2n-1$ and $139$ are odd. Then $2(n+21)-1=139+42(m+1)+i$, and by Theorem~\ref{thm:block} and $P_i=0$ both $a(2n-1)$ and $a(2(n+21)-1)$ equal $Q_i$.

\emph{(iii)} (\decl{even_part}.) For $j<42$ put $\delta_a(j)=P_j+Q_j-3645\,Q_j\in\Z$ (\decl{dA}), and let $b(n)=\delta_a((2(n\bmod 21)+29)\bmod 42)$, so that $b(n+21)=b(n)$. Let $n\ge70$ and write $2n=139+42m+i$ with $m\ge0$ and $0\le i<42$. Since $-139\equiv29\pmod{42}$ and $2n\equiv2(n\bmod21)\pmod{42}$, we have $i=(2(n\bmod21)+29)\bmod42$, so $b(n)=\delta_a(i)$. As $2(n+21)=139+42(m+1)+i$ and $Y_{m+1}=3645Y_m+1$,
\[
  a(2(n+21))=P_i(3645Y_m+1)+Q_i=3645\,(P_iY_m+Q_i)+P_i+Q_i-3645\,Q_i=3645\,a(2n)+b(n).
\]

\emph{(iv)} (\decl{sums_part}.) For $j<42$ put $\delta_s(j)=R_j+U_j-3645\,U_j\in\Z$ (\decl{dS}), and define $c(n)=\delta_s((n+29)\bmod42)$; then $c(n+42)=c(n)$.
For $n\ge139$ write $n=139+42m+i$ with $0\le i<42$; then $i=(n+29)\bmod42$, and exactly as in (iii), $s(n+42)=R_i(3645Y_m+1)+U_i=3645\,s(n)+c(n)$.
For $n=138$: Theorem~\ref{thm:block} with $m=0$ gives $s(139)=U_0=315591$, $a(139)=Q_0=3$ and $s(180)=U_{41}=1150350168$, so $s(138)=315588$; and $c(138)=\delta_s(41)=31908$, and indeed $3645\cdot315588+31908=1150350168=s(180)$. \qed

\begin{remark}[sharpness; not formalized]\label{rem:sharp}
The thresholds are optimal: $a(179)=5\ne101=a(137)$, $a(180)-3645\,a(138)\ne a(222)-3645\,a(180)$ and $s(179)-3645\,s(137)\ne s(221)-3645\,s(179)$, so (ii)--(iv) fail with thresholds $69,69,137$ (checked in Python, also by the math0 referees).
\end{remark}

\section{Machine verification}\label{sec:lean}

\subsection*{Files and declarations}
The proof is in \path{lean/Math0/Frontier/OEIS/A351743.lean} of the math0 repository, \url{https://github.com/Foundation-0/math0}, which imports the frozen statement file (Section~\ref{sec:statement}) and the audit commands in \path{lean/Math0/Audit.lean}. The main declaration is \decl{Math0.Frontier.A351743.result} of type \decl{Math0.Statements.A351743.statement}. The auxiliary declarations of this file named in Section~\ref{sec:proof}, including Theorem~\ref{thm:block}, are \texttt{private}. Kernel computations (\decl{good_all}, \decl{state_base} by \texttt{decide +kernel}; small \texttt{decide}/\texttt{rfl} steps) add no axioms; there is no \texttt{native\_decide}, \texttt{sorry}, \texttt{axiom} or \texttt{set\_option}.

\subsection*{Audit}
The file ends with the command
\begin{lstlisting}
#math0_settles Math0.Frontier.A351743.result Math0.Statements.A351743.statement
\end{lstlisting}
which is used for statements that math0 formalizes itself. It fails unless (1) \decl{statement} is a universe-monomorphic definition of type \texttt{Prop}; (2) \decl{statement} and \decl{result} use no axioms besides \decl{propext}, \decl{Classical.choice} and \decl{Quot.sound} (rejecting \decl{sorryAx}, \decl{Lean.ofReduceBool} from \texttt{native\_decide}, and custom axioms); and (3) the type of \decl{result} is syntactically the statement or its negation (up to metadata, no unfolding). It reports
\begin{lstlisting}[basicstyle=\ttfamily\scriptsize]
math0 audit: 'Math0.Statements.A351743.statement' depends on axioms [Classical.choice, Quot.sound, propext]
math0 audit: 'Math0.Frontier.A351743.result' depends on axioms [Classical.choice, Quot.sound, propext]
math0 statement check: 'Math0.Frontier.A351743.result' proves 'Math0.Statements.A351743.statement'
\end{lstlisting}

\subsection*{Toolchain and reproduction}
Lean~4 \cite{Lean4} \texttt{v4.33.1} and Mathlib \cite{Mathlib} \texttt{v4.33.1}, pinned in \path{lean/}; formal-conjectures \cite{FC} is pinned at commit \texttt{c77fdd4} but not used here.
To reproduce, run \path{scripts/verify.sh} from the repository root; it fetches the Mathlib cache, rebuilds all math0 results with \texttt{lake build} (re-running every audit) and checks that the pinned dependencies are unmodified. This file alone: \texttt{lake env lean Math0/Frontier/OEIS/A351743.lean} in \path{lean} (under two minutes on our machine).

{\footnotesize
\begin{thebibliography}{99}

\bibitem{FC}
M.~Firsching, P.~Lezeau, S.~Mercuri, M.~Z. Horv\'ath, Y.~Dillies, C.~S\"onne, E.~Wieser, F.~Zhang, T.~Hubert, B.~Ag\"uera y Arcas, and P.~Kohli, \emph{Formal Conjectures: An open and evolving benchmark for verified discovery in mathematics}, arXiv:2605.13171, 2026. Repository: \url{https://github.com/google-deepmind/formal-conjectures}.

\bibitem{Fr24}
S.~Fried, \emph{Proofs of some conjectures from the OEIS}, arXiv:2410.07237, 2024.

\bibitem{Fr25}
S.~Fried, \emph{Proofs of several conjectures from the OEIS}, J. Integer Seq.\ \textbf{28} (2025), Article~25.4.3.

\bibitem{Fr26}
S.~Fried, \emph{Proofs of ten conjectures from the OEIS}, J. Integer Seq.\ \textbf{29} (2026), Article~26.1.8.

\bibitem{Fr26b}
S.~Fried, \emph{Further proofs of conjectures from the OEIS}, arXiv:2607.24832, 2026.

\bibitem{Mathlib}
The mathlib Community, \emph{The Lean mathematical library}, in: Proceedings of the 9th ACM SIGPLAN International Conference on Certified Programs and Proofs (CPP 2020), ACM, 2020, 367--381.

\bibitem{Lean4}
L.~de Moura and S.~Ullrich, \emph{The Lean 4 theorem prover and programming language}, in: Automated Deduction -- CADE 28, Lecture Notes in Comput. Sci. \textbf{12699}, Springer, 2021, 625--635.

\bibitem{A351743}
The OEIS Foundation Inc., \url{https://oeis.org/A351743}, entry by D.~J. Sycamore (June~7, 2022), with conjectures by D.~J. Sycamore and R.~Sigrist; revision \#43, June~20, 2022.

\bibitem{A355149}
The OEIS Foundation Inc., \url{https://oeis.org/A355149}, entry by M.~De Vlieger, R.~Sigrist and D.~J. Sycamore (June~21, 2022); revision \#11, Nov.~17, 2022.

\bibitem{A357735}
The OEIS Foundation Inc., \url{https://oeis.org/A357735}, entry by D.~J. Sycamore (Oct.~11, 2022).

\end{thebibliography}
}

\end{document}
